\DocumentMetadata{
  pdfversion=2.0,
  pdfstandard=ua-2,
  lang=en-US,
  tagging=on,
  tagging-setup={math/setup=mathml-SE,tabsorder=structure}
}
\documentclass[11pt,letterpaper]{article}
\usepackage[margin=0.68in,includefoot]{geometry}
\usepackage{newcomputermodern}
\usepackage{amsmath,amscd,mathtools}
\usepackage{tikz,adjustbox}
\providecommand{\square}{\mdlgwhtsquare}
\usepackage[hidelinks]{hyperref}
\hypersetup{
  pdftitle={Analysis PhD Exam, August 1995},
  pdfauthor={University of Florida Department of Mathematics},
  pdfsubject={Graduate mathematics examination},
  pdfdisplaydoctitle=true,
  bookmarksopen=true
}
\setcounter{secnumdepth}{0}
\setlength{\parindent}{0pt}
\setlength{\parskip}{0.28em}
\setlength{\emergencystretch}{3em}
\setlength{\leftmargini}{2.15em}
\setlength{\leftmarginii}{2.65em}
\setlength{\leftmarginiii}{3em}
\pagestyle{empty}
\raggedbottom
\allowdisplaybreaks
\NewTaggingSocketPlug{block/list/label}{examproblem}{%
  \tagstructbegin{tag=Lbl}\tagstructend%
  \tagstructbegin{tag=\LBody}%
  \tagstructbegin{tag=\ProblemHeadingTag,title-o={\ProblemHeadingText}}%
  \tagmcbegin{tag=\ProblemHeadingTag,actualtext={\ProblemHeadingText}}%
  #2%
  \tagmcend%
  \tagstructend%
}
\newenvironment{examproblems}[2]{%
  \def\ProblemHeadingTag{#2}%
  \AssignTaggingSocketPlug{block/list/label}{examproblem}%
  \begin{enumerate}%
  \setcounter{enumi}{#1}%
  \renewcommand{\labelenumi}{\textbf{\theenumi.}}%
  \setlength{\topsep}{0.35em}%
  \setlength{\partopsep}{0pt}%
  \setlength{\itemsep}{0.48em}%
  \setlength{\parsep}{0.18em}%
}{\end{enumerate}}
\newenvironment{examsubparts}{%
  \AssignTaggingSocketPlug{block/list/label}{default}%
  \begin{enumerate}%
  \setlength{\topsep}{0.18em}%
  \setlength{\partopsep}{0pt}%
  \setlength{\itemsep}{0.16em}%
  \setlength{\parsep}{0.1em}%
}{\end{enumerate}}
\newenvironment{examinstructions}{%
  \par\begingroup\itshape
}{%
  \par\endgroup\vspace{0.18em}
}
\makeatletter
\renewcommand\subsection{\@startsection{subsection}{2}{0pt}%
  {-1.25ex plus -.25ex minus -.1ex}%
  {0.55ex plus .1ex}%
  {\normalfont\large\bfseries}}
\renewcommand{\@maketitle}{%
  \newpage\null\vskip -1em%
  {\footnotesize\bfseries
    \href{https://www.ufl.edu/}{University of Florida}, \href{https://math.ufl.edu/}{Department of Mathematics}\par}%
  \vskip -0.5em%
  \tagstructbegin{tag=H1,title={Analysis PhD Exam, August 1995}}%
  \tagpdfparaOff%
  \tagmcbegin{tag=H1}%
    {\LARGE\bfseries\href{https://gma.math.ufl.edu/exams/analysis-phd/}{Analysis PhD Exam}, August 1995\par}%
  \tagmcend%
  \tagpdfparaOn%
  \tagstructend%
  \vskip 0.25em%
  \tagmcbegin{artifact}%
    \hrule height 1pt%
  \tagmcend%
  \vskip 1em%
}
\makeatother
\title{Analysis PhD Exam, August 1995}
\author{}
\date{}
\begin{document}
\maketitle
\thispagestyle{empty}
\begin{examinstructions}
All work must be presented in a neat and logical fashion to receive credit. Be sure to give reasons for all your steps. DO NOT leave any gaps!
\end{examinstructions}
\begin{examinstructions}
\((X,\Sigma,\mu)\) will always denote a measure space (possibly infinite).
\end{examinstructions}
\begin{examproblems}{0}{H2}
\def\ProblemHeadingText{Problem 1.}
\item
Prove that \(L^1(X,\Sigma,\mu)\) is complete.
\def\ProblemHeadingText{Problem 2.}
\item
Let~\(f\) be Lebesgue integrable on \(\mathbb R\). Suppose \(\int_0^x f\,dm=0\) for all real~\(x\), where~\(m\) is Lebesgue measure. What can you say about~\(f\)? Prove your answer.
\def\ProblemHeadingText{Problem 3.}
\item
Let \((f_n)\) be a sequence of integrable functions such that \(\sum_{n=1}^{\infty}\int |f_n|\,d\mu<\infty\). What can you prove about the a.e. \(\mu\) convergence of \(\sum f_n\)?
\def\ProblemHeadingText{Problem 4.}
\item
Let~\(f\) be a Lebesgue integrable function on \(\mathbb R\). Let \(\epsilon>0\). Is it possible to find a bounded interval~\(I\) such that whenever~\(E\) is a measurable set such that \(E\cap I=\phi\), then \(\left|\int_E f\,m\right|<\epsilon\)?
\def\ProblemHeadingText{Problem 5.}
\item
State and prove the Lebesgue Monotone Convergence Theorem.
\def\ProblemHeadingText{Problem 6.}
\item
Present the steps in the construction of the product measure of two finite measure spaces. Sketch the proof of the countable additivity of the product measure.
\def\ProblemHeadingText{Problem 7.}
\item
Let \(X=Y=\) any uncountable set. Let \(\mathcal S=\mathcal T=\) the class of all countable subsets of~\(X\) and~\(Y\), respectively. Note that \(\mathcal S\) and \(\mathcal T\) are~\(\sigma\)-rings. Let \(D=\{(x,y):x=y\}\) be the diagonal in \(X\times Y\). Show that~\(D\) is NOT \(\mathcal S\times\mathcal T\) measurable.
\def\ProblemHeadingText{Problem 8.}
\item
Let \((X,\Sigma,\mu)\) be a finite measure space and \(\mathcal G\) is a sub~\(\sigma\)-algebra of~\(\Sigma\). Define the closed subspace~\(V\) of \(L^\infty(X,\Sigma,\mu)\) as the collection of all \(g\in L^\infty(X,\Sigma,\mu)\) such that \(\int fg\,\mu=0\) for all \(f\in L^1(X,\mathcal G,\mu)\). Prove that if \(k\in L^1(X,\Sigma,\mu)\) and \(\int kg\,d\mu=0\) for all \(g\in V\), then~\(k\) is \(\mathcal G\)-measurable. (Hint: Use the representation for the dual of \(L^1(X,\Sigma,\mu)\) and a corollary of the Hahn–Banach theorem applied to the closed subspace \(L^1(X,\mathcal G,\mu)\) of \(L^1(X,\Sigma,\mu)\).)
\end{examproblems}
\end{document}
